\documentclass[10pt,a4paper]{amsart}
\usepackage[margin=19mm]{geometry}
\usepackage{amsmath,amssymb,mathtools}
\usepackage[scr=rsfs,cal=euler]{mathalfa}
\usepackage{xfrac}
\usepackage[unicode,hidelinks,bookmarks=false]{hyperref}
\newcommand{\ie}{i.e.\ }
\newcommand{\cf}{cf.\ }
\newcommand{\resp}{respectively\ }
\newcommand{\Subsection}{\subsection}
\providecommand{\nobreakdash}{\nobreak\hbox{-}}
\setlength{\emergencystretch}{2em}
\multlinegap0pt
\usepackage{enumerate}
\usepackage[normalem]{ulem}
\usepackage{amssymb}
\newtheorem{thm}{Theorem}
\def\epsilon{\varepsilon}
\def\ez{\epsilon}
\title{Bounded Continuous weak quasiregular mappings that fail to be quasiregular}
\author{Stanislav Hencl, Yi Ru-Ya Zhang}
\date{Source: arXiv:2605.01535v1 (CC BY 4.0)}
\begin{document}
\maketitle

\noindent\emph{Source and licence.} Bounded Continuous weak quasiregular mappings that fail to be quasiregular. Version arXiv:2605.01535v1 (CC BY 4.0), distributed by arXiv under the Creative Commons Attribution 4.0 International licence, http://creativecommons.org/licenses/by/4.0/, as recorded in the arXiv licence metadata for this exact version and shown on the abstract page https://arxiv.org/abs/2605.01535v1. Full licence notice: \texttt{licenses/arxiv-2605.01535-CC-BY-4.0.txt}.\par\smallskip

\noindent\textit{Editorial scope.} Counted unit: the article's thm thm (source0.tex lines 259--264). The article's complete original proof of it is delivered with it (source0.tex lines 367--468, one \texttt{proof} environment of 102 lines, which the article places away from the statement). No mathematics is written, repaired or re-derived: the statement and the proof are the article's own lines, in the article's own notation, and no retained line is substituted or re-flowed.  No further source-local context is needed: every cross-reference of every retained line resolves inside the two delivered spans.  Nothing was omitted from any retained span, so the package carries no omission record.  No retained line contains a citation, so the wrapper carries no bibliography and no bibliography entry.  No retained line contains a figure environment, an image-inclusion command or drawing code, so no image is delivered and the package declares no asset dependency.  Editorial formatting only, in addition: an AMS \texttt{amsart} preamble that restates the article's own declarations and macros used by the retained lines, namely the environments \texttt{thm} on separate counters, together with the standard packages those lines need.  The retained environments print in this document as Theorem 1 in order of appearance, whereas the article prints the same items with the numbering of its own full document.  The complete native source of the article as distributed in the arXiv e-print for this exact version is bundled unchanged as \texttt{sources/199642/source0.tex}.
\begin{thm}\label{main thm}
    When $n\ge3$, for any bounded domain $\Omega\subset \mathbb R^n$ and every $p\in \left(1,\,\frac{nK}{K+1}\right)$, there is a weakly $K$-quasiregular mapping 
    $$f\in W^{1,\,p}(\Omega;\,\mathbb R^n)\cap C(\Omega;\,\mathbb R^n)
   \cap L^\infty(\Omega;\mathbb R^n)$$
    which is not quasiregular.
\end{thm}

\begin{proof}[Proof of Theorem~\ref{main thm}]
Given a ball $B(y,r)$, $\mathbf z\in \mathbb R^n$ and $t>0$ we define a mapping
$$\Phi_{B}^{\mathbf z,t}(\mathbf x)=\mathbf z+ t(\mathbf x-\mathbf y)\left(\frac{r}{|\mathbf x-\mathbf y|}\right)^{1+\frac 1 K}.$$
\noindent{\bf Step 1: Construction of  a weakly quasiregular mapping $F$.} Fix $0<a<1$ and consider an exact packing $\mathcal F_1=\{B_{1,\,j}\}_{j=1}^\infty, B_{1,\,j}=B(\mathbf {y}_{1,\,j},\,r_{1,\,j})$ of $\Omega=:U_1$ by balls  with radius no more than $\delta_1>0$. We define 
$$
F_1(\mathbf {x}) = \left\{ \begin{array}{ll}
\Phi^{\mathbf{y}_{1,j},1}_{B_{1,\,j}}(\mathbf {x}) & \textrm{ if } \ \mathbf x\in B_{1,\,j}\setminus a B_{1,\,j}\\
\mathbf {y}_{1,\,j}+a^{-\frac 1 K-1}(\mathbf x-\mathbf {y}_{1,\,j}) & \textrm{ if } \ \mathbf  x\in aB_{1,\,j}\\
\mathbf {x}  & \textrm{ otherwise }
\end{array} \right. .$$
This give us the first continuous map $F_1$.

We next construct $F_2$. Towards this, we consider  an exact packing $\mathcal F_2=\{B_{2,\,j}\}_{j=1}^\infty$ of 
$$U_2=\bigcup_{j=1}^\infty aB_{1,\,j}$$
by balls $B_{2,\,j}=B(\mathbf {y}_{2,\,j},\,r_{2,\,j})$ with radius no more than $\delta_2>0$.
Then we define
$$
F_2(\mathbf {x}) = \left\{ \begin{array}{ll}
\Phi^{F_1(\mathbf y_{2,j}),a^{-\frac{1}{K}-1}}_{B_{2,\,j}}(\mathbf {x}) & \textrm{ if } \ \mathbf x\in B_{2,\,j}\setminus a B_{2,\,j}\\
F_1(\mathbf {y}_{2,\,j})+a^{-\frac 2 K-2}(\mathbf x-\mathbf {y}_{2,\,j}) & \textrm{ if } \  \mathbf x\in aB_{2,\,j}\\
F_1(\mathbf x) & \textrm{ otherwise }
\end{array} \right. .$$
Note that for $\mathbf{x}\in \partial a B_{2,\,j}$ we have
$$
\Phi^{F_1(\mathbf{y}_{2,j}),a^{-\frac{1}{K}-1}}_{B_{2,\,j}}(\mathbf {x})=
F_1(\mathbf{y}_{2,j})+a^{-\frac{1}{K}-1} (\mathbf{x}-\mathbf{y}_{2,j})
\left(\frac{r_{2,\,j}}{ar_{2,\,j}}\right)^{1+\frac 1 K}=F_1(\mathbf {y}_{2,\,j})+a^{-\frac 2 K-2}(\mathbf x-\mathbf {y}_{2,\,j})
$$
and for $\mathbf{x}\in \partial B_{2,\,j}$, since $F_1$ is linear on $U_2$ we have
$$
\Phi^{F_1(\mathbf{y}_{2,j}),a^{-\frac{1}{K}-1}}_{B_{2,\,j}}(\mathbf {x})=
F_1(\mathbf{y}_{2,j})+a^{-\frac{1}{K}-1} (\mathbf{x}-\mathbf{y}_{2,j})
\left(\frac{r_{2,\,j}}{r_{2,\,j}}\right)^{1+\frac 1 K}= F_1(\mathbf{x})
$$
so the mapping is continuous.

Now we iterate this construction. Namely, given $\mathcal F_{k-1}$,   set 
$$U_k=\bigcup_{j=1}^\infty aB_{k-1,\,j},$$
and consider an exact packing $\mathcal F_k=\{B_{k,\,j}\}_{j=1}^\infty$ of
$U_k$ by balls $B_{k,\,j}=B(\mathbf {y}_{k,\,j},\,r_{k,\,j})$ with radius no more than $\delta_k>0$. Then define inductively
$$
F_k(\mathbf {x}) = \left\{ \begin{array}{ll}
\Phi^{F_{k-1}(\mathbf{y}_{k,j}),a^{(k-1)\left(-\frac{1}{K}-1\right)}}_{B_{k,\,j}}(\mathbf {x}) & \textrm{ if } \ \mathbf x\in B_{k,\,j}\setminus a B_{k,\,j}\\
F_{k-1}(\mathbf{y}_{k,j})+a^{-\frac {k} K-k}(\mathbf x-\mathbf {y}_{k,\,j})   & \textrm{ if } \  \mathbf x\in aB_{k,\,j}\\
F_{k-1}(\mathbf x) & \textrm{ otherwise }
\end{array} \right. $$
and again it is not difficult to check that this mapping is continuous. 
We further require that $\delta_k$ are chosen to satisfy
\begin{equation}\label{finite sum}
    \sum_{k=1}^\infty a^{-\frac k K-k} \delta_k<\infty. 
\end{equation}
This gives us a sequence of continuous mappings  $\{F_k\}$.


Moreover,  we claim that $F_k$ are converging uniformly. Indeed, according to the definition of $F_k$, the image of 
\begin{equation}\label{image}
F_k(B_{k,\,j}\setminus a B_{k,\,j})\subset  B(F_{k-1}(\mathbf{y}_{k,j}),\,a^{-\frac k K-k}\delta_k),
\end{equation}
and by \eqref{finite sum} it it easy to see that $F_k$ are uniformly bounded. Using the  Cauchy convergence \eqref{finite sum} it is also standard to show that $F_k$ converge uniformly to a continuous function $F$. 

%In addition, given any $\ez>0$, let $k_\ez>0$ so that
%$$\sum_{k=k_\ez}^\infty a^{-\frac k K-k} \delta_k\le \ez.$$
%Then by choosing $d_\ez=\max_{k\ge k_\ez}\{\delta_k\},$ it follows from the definition %of $F_k$ that, for any pair of points satisfying $\mathbf{x}_1,\,\mathbf{x}_2 \in \Omega,%\,|\mathbf{x}_1-\mathbf{x}_2|\le d_\ez$, one has
%$$|d(\mathbf{x}_1)-d(\mathbf{x}_2)|\le \sum_{k=k_\ez}^\infty a^{-\frac {k} K-k} %%\delta_k\le \ez.$$
%Namely, $F_k$ are equicontinuous, and then we concludes the claim. 
%Therefore, up to passing to a subsequence, $F_k$ converges to a continuous function $F$ %locally uniformly. 

In addition, according to our construction, it follows that
\begin{equation}\label{zero null set}
    \left|\Omega\setminus \bigcup_{k=1}^\infty\bigcup_{j=1}^\infty (B_{k,\,j}\setminus aB_{k,\,j})\right|=0,
\end{equation} 
and for every $\mathbf x\in B_{k,\,j}\setminus aB_{k,\,j}$,
\begin{equation}\label{easy}
F_{l}(\mathbf x) = F_k(\mathbf x)= F_{k-1}(\mathbf{y}_{k,\,j})+  a^{-(k-1) (1+\frac 1 K)}(\mathbf x-\mathbf y_{k,\,j})\left(\frac{r_{k,\,j}}{|\mathbf x-\mathbf y_{k,\,j}|}\right)^{1+\frac 1 K}, \quad \forall l\ge k,
\end{equation}
which coincides with $\Phi_{B_{k,\,j}}$ up to a translation and a dilation. 
Since each $\Phi_{B_{k,\,j}}$ is $K$-quasiregular and $\det DF\le 0$ almost everywhere, it follows that $F$ is weakly $K$-quasiregular.  However, $F$ is not $K$-quasiregular as its degree is positive on balls $B_{k,j}$ as our $F$ is just a positive linear transformation (i.e. $F(x)=a_{k,j}+a^{-(k-1)(\frac{1}{K}+1)} x$, see \eqref{easy}) on $\partial B_{k,j}$. 

\medskip

\noindent{\bf Step 2: Energy estimate of $F$.} Now we compute $|DF|$. Observe that when $\mathbf x\in B_{k,\,j}\setminus a B_{k,\,j}$
$$D\Phi_{B_{k,\,j}}(\mathbf x) =a^{-\frac {K+1} K(k-1)}\left(\frac{r_{k,\,j}}{|\mathbf {x}-\mathbf {y}_{k,\,j}|}\right)^{1+\frac 1 K} \left[Id-\frac{K+1}{K} \frac{(\mathbf {x}-\mathbf {y}_{k,\,j})\otimes (\mathbf {x}-\mathbf {y}_{k,\,j})}{|\mathbf {x}-\mathbf {y}_{k,\,j}|^2}\right].$$
Thus
$$|D\Phi_{B_{k,\,j}}(\mathbf x)|\le C(K)a^{-\frac {K+1} K(k-1)} r_{k,\,j}^{1+\frac 1 K}|\mathbf {x}-\mathbf {y}_{k,\,j}|^{-\left(1+\frac 1 K\right)},$$
and then
\begin{multline}\label{energy in one}
\int_{B_{k,\,j}\setminus a B_{k,\,j}} |D\Phi_{B_{k,\,j}}|^p\,dx\le  C(n,\,K) a^{-\frac {K+1} K(k-1)p} r_{k,\,j}^{p\left(1+\frac 1 K\right)}\int_{ar_{k,\,j}}^{r_{k,\,j}} s^{-p\left(1+\frac 1 K\right)+n-1}\, ds \\
\le C(n,\,K,\,a,\,p) a^{-\frac {K+1} K kp} r_{k,\,j}^{p\left(1+\frac 1 K\right)} r_{k,\,j}^{-p\left(1+\frac 1 K\right)+n}(1-a^{-p\left(1+\frac 1 K\right)+n})\le C(n,\,K,\,a,\,p) a^{-\frac {K+1} K kp}r_{k,\,j}^n,
\end{multline}
where we applied
$$-p\left(1+\frac 1 K\right)+n-1> -1 \quad \Longleftrightarrow \quad p<\frac{nK}{K+1}.$$
Now by summing over all $B_{k,\,j}$ and applying \eqref{energy in one} together with \eqref{zero null set}, we conclude that, for each $k_0\in \mathbb N$,
\begin{align*}
\int_{\Omega} |DF_{k_0}|^p\,dx\le & \ \sum_{k=1}^{k_0}\sum_{j=1}^\infty\int_{B_{k,\,j}\setminus a B_{k,\,j}} |D\Phi_{B_{k,\,j}}|^p\,dx + \sum_{j=1}^{\infty}\int_{B_{k_0,\,j}} |D\Phi_{B_{k_0,\,j}}|^p\,dx \\
\le  & \ C(n,\,K,\,a,\,p)\sum_{k,\,j=1}^\infty a^{-kp\frac{K+1} K }r_{k,\,j}^n  \\ 
= & \ C(n,\,K,\,a,\,p)\sum_{k=1}^\infty a^{-kp\frac{K+1} K }|U_{k}|\\
=&\ C(n,\,K,\,a,\,p)\sum_{k=1}^\infty a^{-kp\frac{K+1} K } a^{nk}|U_1|
\le C(n,\,K,\,a,\,p) |\Omega|<\infty, 
\end{align*}
where we applied that $\mathcal F_k$ is an exact packing of $U_k$ and $\frac{K}{K+1}n>p$.
Thus, we conclude $F\in W^{1,\,p}(\Omega)$ by the weak convergence of $F_k\rightharpoonup F$ in $W^{1,\,p}(\Omega)$, up to further passing to a subsequence. 
\end{proof}

\end{document}
